\section{Compilation Time}\label{sec:compile-time}

Table~\ref{tab:benchmark-results} shows the measured results.
The fastest case and the realistic example are close to 2 ms.
The largest synthetic cases are around 6 ms.
Because these times are small, they should be interpreted as local measurements on the stated machine rather than as a
portable performance result.
They are still useful for the thesis because they show that the current native compiler is fast enough to support an
interactive compile loop in this local setup.

\begin{table}[htbp]
  \centering
  \footnotesize
  \renewcommand{\arraystretch}{1.25}
  \begin{tabular}{lrrrrr}
    \toprule
    \textbf{Program} & \textbf{LOC} & \textbf{Note events} & \textbf{MIDI size} & \textbf{Mean ms} & \textbf{SD ms} \\
    \midrule
    \texttt{single\_loop\_100} & 7 & 100 & 856 B & 2.02 & 0.31 \\
    \texttt{single\_loop\_1000} & 7 & 1,000 & 8,056 B & 2.14 & 0.31 \\
    \texttt{single\_loop\_10000} & 7 & 10,000 & 80,056 B & 3.96 & 0.24 \\
    \texttt{single\_loop\_20000} & 7 & 20,000 & 160,056 B & 6.03 & 0.15 \\
    \texttt{chord\_loop\_5000} & 7 & 20,000 & 160,056 B & 5.90 & 0.17 \\
    \texttt{drum\_groove\_500} & 18 & 10,000 & 88,053 B & 6.39 & 0.19 \\
    \texttt{example\_copy} & 77 & 188 & 1,751 B & 2.04 & 0.18 \\
    \bottomrule
  \end{tabular}
  \caption{Local Release-build compilation benchmark.
  Times include MIDI file writing and are based on 50 measured runs after 5 warmups.}\label{tab:benchmark-results}
\end{table}

\begin{figure}[htbp]
  \centering
  \begin{tikzpicture}
    \begin{axis}[
        ybar,
        width=0.95\textwidth,
        height=7cm,
        ymin=0,
        ylabel={Mean compile time (ms)},
        symbolic x coords={100,1000,10000,20000,chord,drums,example},
        xtick=data,
        x tick label style={rotate=35, anchor=east, font=\scriptsize},
        nodes near coords,
        nodes near coords style={font=\scriptsize},
        bar width=12pt,
        enlarge x limits=0.08,
        ymajorgrids=true,
        grid style={draw=motivoInk!12}
      ]
      \addplot[fill=motivoBlue!70, draw=motivoBlue!90] coordinates {
        (100,2.02)
        (1000,2.14)
        (10000,3.96)
        (20000,6.03)
        (chord,5.90)
        (drums,6.39)
        (example,2.04)
      };
    \end{axis}
  \end{tikzpicture}
  \caption{Mean compile times for the local benchmark corpus.}\label{fig:benchmark-time}
\end{figure}

The near-linear growth in the single-loop cases is consistent with the implementation.
Parsing and semantic analysis see a small source program, while lowering and MIDI writing scale with the number of
generated events.
The backend also sorts events by tick before computing delta-times.
The result is still small at the current safety limit, but the implementation should not be described as unbounded.

\subsection{Execution Speed and Real-Time Viability}\label{detail-subsec:execution-speed-and-real-time-viability}

Because the compiler is implemented as a native C++23 program without reliance on garbage collection or dynamic dispatch
for AST traversal, the constant factors in the time complexity are small in the current measurements.
In practice, the compiler generates compositions containing event streams up to the current implementation
limit in a fraction of a second.

These local measurements are relevant to Motivo Studio because the editor workflow depends on compilation
remaining quick enough for interactive use.
The interactive playground environment (discussed in previous chapters) relies on the compiler's ability to re-compile
the entire source file during interactive compile actions.
Because the execution time for typical musical pieces (ranging from hundreds to thousands of events) sits
well below the latency usually noticed during ordinary editor interaction, the composer experiences
interactive auditory and visual feedback without using an interpreted runtime during playback.
